%% sections/08_results.tex

\section{측정 결과}
\label{sec:results}

전 결과는 DGX B200$\times$8 실측이다 (\S\ref{sec:methodology}).
주력 8모델 $\times$ 4방법 $\times$ 7조건 + XL 2모델 $\times$ 2방법 $\times$ 7조건.

%% -----------------------------------------------------------------------
\subsection{개요: 55 net-positive / 15 vanilla-win}
\label{subsec:res-overview}

70개 (모델$\times$조건) 쌍에서 suffix 는 55개 net-positive($+4\!\sim\!+232\%$),
15개에서 vanilla 가 더 빠르다(\S\ref{subsec:mot-heterogeneous}, Fig.~\ref{fig:alpha-gain}).
mix 조건만 보면 8개 주력 모델 전부 net-positive 이나(Table~\ref{tbl:b200-realtrace};
XL 중 DeepSeek-R1 671B 은 mix 도 음, \S\ref{subsec:res-modeltype}),
단일 corpus 로 내려가면 부호가 갈린다.
이 절은 그 부호를 \emph{예측}하고(\S\ref{subsec:res-alpha}),
실패 모델을 \emph{식별}하며(\S\ref{subsec:res-modeltype}),
성공이 남기는 \emph{자원 여유}를 정량화한다(\S\ref{subsec:res-resource}).

%% -----------------------------------------------------------------------
\subsection{예측 조건: $\alpha$ 대 overhead}
\label{subsec:res-alpha}

%% tables/tbl_b200_alpha.tex  (Table IV) — width-safe: 수치 2행만 표, 15셀 목록은 caption
%% source: 70 (model×condition) 쌍

\begin{table}[t]
\centering
\caption{\textbf{언제 spec-decode가 이득인가}: 70개 (모델$\times$조건) 쌍을
  suffix net-positive / vanilla-win 으로 분리. 부호는 accept rate $\alpha$ 와
  per-step overhead(wall ratio)의 경쟁으로 결정된다.
  vanilla-win 15셀은 모두 추론/MoE 모델이다 ---
  R1-671B(MoE) swebench$-63.5\%$/mix$-49.2\%$($\alpha0.36$--$0.54$, 7/7),
  DS-R1-Distill-L70B swebench$-15.4$/humaneval$-2.2$($\alpha0.27$--$0.38$, 6/7),
  Q72B wildchat$-6.5$/mbpp$-4.8$.}
\label{tbl:b200-alpha}
\resizebox{\columnwidth}{!}{%
\begin{tabular}{@{}lrcc@{}}
\toprule
\textbf{그룹} & \textbf{셀 수} & $\boldsymbol{\alpha}$ \textbf{(중앙 / 범위)} & \textbf{wall$_s$/wall$_v$} \\
\midrule
suffix net-positive & 55 & $0.68$ \;($0.28$--$1.41$) & $0.60$ \\
vanilla wins         & 15 & $0.37$ \;($0.27$--$0.54$) & $1.18$ \;($\le 2.9$) \\
\bottomrule
\end{tabular}}
\end{table}


Table~\ref{tbl:b200-alpha} 와 같이 net-positive 셀의 수락률 $\alpha$ 중앙값은 $0.68$
($0.28\!\sim\!1.41$), vanilla-win 셀은 $0.37$($0.27\!\sim\!0.54$)이다.
그러나 break-even 은 단일 $\alpha$ 값이 아니다:
net-positive 셀의 wall ratio(suffix/vanilla) 중앙값은 $0.60$,
vanilla-win 셀은 $1.18$(최대 $2.9$)로, \emph{overhead 가 큰 모델일수록 더 높은 $\alpha$ 를 요구}한다
(식~\ref{eq:rk-condition}).
실제로 한 저-overhead 셀은 $\alpha\!=\!0.28$ 에서도 $+209\%$ 인 반면,
DeepSeek-R1 671B 는 $\alpha\!\approx\!0.45\!\sim\!0.54$ 에서도 손해다.
즉 $\Delta(\mathrm{suffix})$ 의 부호는 $\alpha\,K_{\mathrm{eff}}$ 대 overhead 의 경쟁으로 설명된다.

%% -----------------------------------------------------------------------
\subsection{모델 유형이 부호를 결정한다}
\label{subsec:res-modeltype}

%% tables/tbl_b200_xl.tex  (Table V) — XL: 405B FP8(dense) vs DeepSeek-R1 671B(MoE 추론)
%% source: TSK_042 XL 매트릭스

\begin{table}[t]
\centering
\caption{XL 스케일 (vanilla/suffix, tps).
  동일 $\sim\!400\text{--}700$B 급에서 \emph{dense 표준}(Llama-3.1-405B-FP8)은 suffix 가 전 조건 net-positive,
  \emph{MoE 추론}(DeepSeek-R1)은 전 조건 net-negative.
  최적 방법은 파라미터 수가 아니라 \emph{모델 유형}이 결정한다.
  두 모델 모두 CPU 이용률 $\le 4.8\%$.
  llm-d 도 7조건 측정했으며 그 한계는 \S\ref{subsec:res-llmd} 에서 다룬다.}
\label{tbl:b200-xl}
\footnotesize
\setlength{\tabcolsep}{5pt}
\resizebox{\columnwidth}{!}{%
\begin{tabular}{@{}lrrrr@{}}
\toprule
\textbf{조건} & \multicolumn{2}{c}{\textbf{Llama-3.1-405B-FP8 (dense)}} & \multicolumn{2}{c}{\textbf{DeepSeek-R1 671B (MoE)}} \\
 & van$\to$suf & $\Delta$ & van$\to$suf & $\Delta$ \\
\midrule
sharegpt  & $1217\to2061$ & $+69\%$  & $1475\to797$ & $-46\%$ \\
wildchat  & $1280\to2290$ & $+79\%$  & $1556\to858$ & $-45\%$ \\
lmsys     & $1220\to2243$ & $+84\%$  & $1533\to811$ & $-47\%$ \\
swebench  & $1204\to2639$ & $+119\%$ & $1474\to538$ & $-64\%$ \\
mbpp      & $\;916\to1725$ & $+88\%$ & $1437\to677$ & $-53\%$ \\
humaneval & $1253\to2112$ & $+69\%$  & $1004\to606$ & $-40\%$ \\
mix       & $1252\to2829$ & $+126\%$ & $1538\to781$ & $-49\%$ \\
\midrule
$\alpha$ (mix) & \multicolumn{2}{c}{$0.766$} & \multicolumn{2}{c}{$0.451$} \\
\bottomrule
\end{tabular}}
\end{table}


vanilla-win 15셀은 \emph{모두} 추론 특화 또는 MoE 모델이다
(DeepSeek-R1 671B 7/7, DeepSeek-R1-Distill-Llama-70B 6/7, Qwen-2.5-72B 2셀; Table~\ref{tbl:b200-alpha}).
XL 스케일이 이를 가장 선명하게 보인다(Table~\ref{tbl:b200-xl}):
동일 $\sim\!400\!\sim\!700$B 급에서 \emph{dense 표준} Llama-3.1-405B-FP8 은 suffix 가
전 조건 $+69\!\sim\!+126\%$ 인 반면, \emph{MoE 추론} DeepSeek-R1 671B 는
전 조건 $-40\!\sim\!-64\%$ 다.
파라미터 수가 더 큰 R1 이 더 나쁘므로, \textbf{최적 방법은 크기가 아니라 모델 유형이 결정한다}.
이 (모델$\times$조건) 의존성은 단일 고정 방법이 구조적으로 최적일 수 없음을 입증하며
(Property~\ref{prop:no-dominant}), 게이트의 1차 입력이 모델 유형임을 함의한다.
실패 메커니즘은 \S\ref{sec:mechanism} 에서 분석한다.

%% -----------------------------------------------------------------------
\subsection{llm-d 라우팅: 복잡성의 한계 이득}
\label{subsec:res-llmd}

vanilla$+$suffix 두 백엔드를 동일 GPU 예산으로 분할한 llm-d(\S\ref{sec:methodology})는
전 10모델 70개 셀에서 vanilla 대비 51개(73\%)에서 우위이나,
\emph{최선 단일 방법}을 능가하는 경우는 10개(14\%)에 그치며
주로 7B$\cdot$32B 의 대화 조건에 국한된다.
복잡한 클러스터 라우팅의 추가 이득은 제한적이며,
서버 내 경량 모델-유형 인지 선택이 대부분의 가치를 포착한다(\S\ref{sec:discussion}).
특히 DeepSeek-R1-Distill-Llama-70B 처럼 한 백엔드(suffix)가 열등하면
라우터 평균도 vanilla 단일보다 낮아진다(전 조건 vanilla 우위).

XL 스케일에서도 동일 패턴이 확인된다(mix): Llama-3.1-405B-FP8 는 llm-d $1{,}429$ 가
suffix $2{,}829$ 에 크게 못 미치고, DeepSeek-R1 671B 는 llm-d $1{,}008$ 이
suffix $781$ 은 능가하나 \emph{vanilla $1{,}538$ 에는 미달}한다.
즉 지배 방법이 한쪽으로 분명할 때 두-백엔드 라우터는 최선 단일을 넘지 못한다.

%% -----------------------------------------------------------------------
\subsection{llm-d 대체: 모델-유형 게이트 (측정 예정)}
\label{subsec:res-gate}

llm-d 의 한계 이득(14\%)은 무거운 disaggregation 대신 경량 게이트로 대체 가능함을 시사한다
(\S\ref{subsec:gate}).
제안 게이트의 처리량 \emph{상한}(oracle$=$best-single, 본 연구 실측)은
이미 전 모델에서 llm-d 를 능가한다 (Table~\ref{tbl:gate}):
mix 기준 Llama-3.1-70B $+160\%$, DeepSeek-R1-Distill-Llama-70B $+114\%$,
DeepSeek-R1 671B $+53\%$, Qwen-2.5-32B $+26\%$ 로,
지배 방법이 분명한 모델일수록 격차가 크다.
실분류기를 적용한 \emph{실현} 게이트 처리량과 oracle 대비 decision-regret 은
SUB\_200 (TSK\_044 CPU 분류기) 후속 측정 대상이며, 본 연구는 222 셀 oracle
(Table~\ref{tbl:gate}) 을 분류기 학습 데이터로 확정한다.
분류 오버헤드는 라우터 host 측에서 ms/req 단위로 측정 예정이며, 본 oracle 의
\emph{상한}만으로도 llm-d 대비 Llama-3.1-70B $+160\%$, DeepSeek-R1 671B $+53\%$
의 격차가 분명하다.

%% -----------------------------------------------------------------------
\subsection{자원 signature: 성공이 GPU를 비우고, CPU는 늘 유휴}
\label{subsec:res-resource}

%% tables/tbl_b200_resource.tex
%% 자원 signature: spec-decode 성공 = GPU un-saturation, 실패 = GPU saturated. CPU 전구간 idle.
%% source: 70 (model×condition) 쌍 GPU util(vanilla->suffix) + 140셀 CPU util

\begin{table}[t]
\centering
\caption{자원 signature.
  spec-decode 가 \emph{성공}하면 GPU 가 un-saturate(여유 발생)되고,
  \emph{실패}하면 GPU 는 거부된 draft 의 낭비 연산으로 포화 상태를 유지한다.
  CPU 는 7B부터 671B까지 전 구성에서 유휴($\le 4.9\%$).}
\label{tbl:b200-resource}
\footnotesize
\setlength{\tabcolsep}{5pt}
\begin{tabular}{@{}lrr@{}}
\toprule
\textbf{그룹} & \textbf{GPU\% (van$\to$suf, 중앙)} & \textbf{함의} \\
\midrule
suffix net-positive (55) & $94.9 \to 78.7$ ($-16$pp) & GPU slack 발생 \\
vanilla wins (15)        & $98.2 \to 86.8$ ($-11$pp) & 낭비 연산으로 포화 \\
\midrule
CPU util (전 140셀) & \multicolumn{2}{r}{중앙 $4.4\%$, 범위 $2.5\text{--}4.9\%$ (XL 포함)} \\
\bottomrule
\end{tabular}
\end{table}

%% tables/tbl_b200_oracle.tex  (조건별 oracle, 전 10모델 70셀)
%% source: tput + routing_combined + routing_llmd merge (canonical, 2026-06-04)

\begin{table*}[t]
\centering
\caption{조건별 oracle 방법 (10모델 $\times$ 7조건 $=$ 70셀, 처리량 최대 방법).
  요청 출처와 모델 유형에 따라 최적이 달라져 단일 고정 방법은 최적일 수 없다.
  \textbf{V}=vanilla, \textbf{S}=suffix, \textbf{L}=llm-d.
  DeepSeek-R1(MoE 추론)은 전 조건 vanilla, DeepSeek-R1-Distill-Llama-70B 은 mix 외 vanilla;
  dense 표준은 대부분 suffix, 중형 Qwen 대화 조건은 llm-d.}
\label{tbl:b200-oracle}
\scriptsize
\setlength{\tabcolsep}{5pt}
\begin{tabularx}{\textwidth}{@{}Xccccccc@{}}
\toprule
\textbf{Model} & \textbf{ShareGPT} & \textbf{WildChat} & \textbf{LMSYS} &
  \textbf{SWE-bench} & \textbf{MBPP} & \textbf{HumanEval} & \textbf{mix} \\
\midrule
Qwen-2.5-7B-Instruct          & L & L & L & L & S & S & S \\
Llama-3.1-8B-Instruct         & S & S & S & S & S & S & S \\
DeepSeek-R1-Distill-Qwen-7B   & S & S & L & S & S & S & S \\
Qwen-2.5-32B-Instruct         & L & L & L & S & L & S & S \\
DeepSeek-R1-Distill-Qwen-32B  & S & S & S & S & S & S & S \\
Qwen-2.5-72B-Instruct         & S & V & S & S & L & S & S \\
Llama-3.1-70B-Instruct        & S & S & S & S & S & S & S \\
DeepSeek-R1-Distill-Llama-70B & V & V & V & V & V & V & S \\
Llama-3.1-405B-Instruct-FP8   & S & S & S & S & S & S & S \\
DeepSeek-R1 (671B, MoE)       & V & V & V & V & V & V & V \\
\bottomrule
\end{tabularx}
\end{table*}

%% tables/tbl_b200_latency.tex  (Table VIII) — width-safe: \resizebox
%% source: TSK_042 routing_combined (mix)

\begin{table}[t]
\centering
\caption{실 트레이스(mix) 지연 (p50, ms)와 suffix accept $\alpha$.
  suffix 는 드래프팅으로 TTFT 가 소폭 증가하나 다중 토큰 수락으로 TPOT 를 크게 단축.}
\label{tbl:b200-latency}
\resizebox{\columnwidth}{!}{%
\begin{tabular}{@{}lccc@{}}
\toprule
\textbf{Model} & \textbf{TTFT (v$\to$s)} & \textbf{TPOT (v$\to$s)} & \textbf{$\alpha_{\text{suf}}$} \\
\midrule
Llama-3.1-8B-Instruct          & $22.8\to24.7$ & $3.5\to1.0$ & $0.933$ \\
Llama-3.1-70B-Instruct         & $28.4\to56.3$ & $9.2\to2.5$ & $0.915$ \\
DeepSeek-R1-Distill-Qwen-7B    & $16.8\to22.3$ & $3.3\to0.9$ & $0.876$ \\
DeepSeek-R1-Distill-Qwen-32B   & $24.1\to37.0$ & $5.9\to1.6$ & $0.801$ \\
DeepSeek-R1-Distill-Llama-70B  & $28.4\to50.1$ & $9.0\to2.1$ & $0.786$ \\
Qwen-2.5-7B-Instruct           & $26.2\to69.3$ & $6.9\to3.1$ & $0.881$ \\
Qwen-2.5-32B-Instruct          & $29.6\to65.4$ & $9.3\to3.1$ & $0.857$ \\
Qwen-2.5-72B-Instruct          & $31.2\to57.4$ & $9.3\to2.6$ & $0.852$ \\
\bottomrule
\end{tabular}}
\end{table}


spec-decode 의 성공/실패는 자원 균형에 상반된 흔적을 남긴다(Table~\ref{tbl:b200-resource}).
\textbf{성공 셀}에서 GPU 이용률은 $94.9\%\!\to\!78.7\%$(중앙, $-16$pp)로 떨어진다 ---
step 수가 줄며 GPU 가 더 자주 쉰다.
\textbf{실패 셀}에서는 $98.2\%\!\to\!86.8\%$ 로 높게 유지된다 ---
거부되는 draft 의 낭비 연산으로 GPU 가 계속 바쁘다.
지연 측면에서도 성공 셀은 TPOT(토큰당 지연)를 크게 줄인다
(예: DeepSeek-R1-Distill-Qwen-7B $3.3\!\to\!0.9$\,ms; Table~\ref{tbl:b200-latency}).

한편 \textbf{CPU 는 7B부터 671B까지 전 140셀에서 $2.5\!\sim\!4.9\%$(중앙 $4.4\%$)} 로 유휴하다.
즉 처리량 지배 기법이 성공할 때 GPU 여유와 CPU 유휴가 \emph{공존}하며,
이 둘이 host-side reclamation 의 회수 대상이 된다(\S\ref{sec:direction}).
(모델, 조건) oracle 전체는 Table~\ref{tbl:b200-oracle} 에 정리한다.

%% -----------------------------------------------------------------------
\subsection{Host-side CPU reclamation: 75+ lever 측정 (모두 measured-negative)}
\label{subsec:res-reclaim}

%% tables/tbl_reclamation.tex  (Table IX) — width-safe: \resizebox + 주석 caption
%% [PLACEHOLDER] Host-side CPU 병렬 reclamation 성능. 결정적 실험(§6.5) 후 채움.

\begin{table}[t]
\centering
\caption{\textbf{[측정 예정]} Host-side CPU 병렬 reclamation 성능.
  suffix 가 un-saturate 시킨 GPU(\S\ref{subsec:res-resource})를
  유휴 CPU 의 병렬 host-path 로 재충전했을 때의 처리량과 GPU 이용률 회복.
  suffix tps$\cdot$GPU\% 는 본 연구 실측(B200); +reclaim 열은
  \S\ref{subsec:decisive} 결정적 실험 후 측정(\textbf{TBD}).}
\label{tbl:reclamation}
\resizebox{\columnwidth}{!}{%
\begin{tabular}{@{}lrccc@{}}
\toprule
\textbf{Model (mix)} & \textbf{suffix tps} & \textbf{+reclaim} &
  \textbf{$\Delta$} & \textbf{GPU\% suf$\to$rcl} \\
\midrule
Qwen-2.5-7B-Instruct   & 7{,}803  & \textit{TBD} & \textit{TBD} & $26.5 \to \textit{TBD}$ \\
Qwen-2.5-32B-Instruct  & 6{,}597  & \textit{TBD} & \textit{TBD} & $64.8 \to \textit{TBD}$ \\
Llama-3.1-8B-Instruct  & 27{,}851 & \textit{TBD} & \textit{TBD} & $62.8 \to \textit{TBD}$ \\
Llama-3.1-70B-Instruct & 10{,}400 & \textit{TBD} & \textit{TBD} & $83.4 \to \textit{TBD}$ \\
\midrule
\textbf{기하평균} & --- & \textit{TBD} & \textit{TBD} & --- \\
\bottomrule
\end{tabular}}
\end{table}


\S\ref{subsec:res-resource} 가 정량화한 자원 여유 ---
spec-decode 성공 셀의 GPU 여유($u_{\mathrm{gpu}}^{\mathrm{suf}}\!=\!26\text{--}65\%$)와
전 구성 유휴 CPU($\le4.9\%$) --- 는
\S\ref{subsec:parallel-model} 의 CPU 병렬 reclamation 으로 처리량을 회수할 \emph{상한}을 정의한다.
본 연구는 이 상한을 실측하기 위해 6 회 누적 sweep 으로
$\ge\!75$ host-path reclamation lever 를 측정하였다 (Table~\ref{tbl:reclamation}).
범위는 env-flag (jemalloc, MALLOC\_ARENA, PYTORCH\_CUDA\_ALLOC\_CONF, FA3, FlashInfer sampler),
NUMA bind / SMT pinning, AVX-512 host helper (BPE encode, simdjson 요청 파싱, base64 출력 streaming),
KV layout (SoA paged, huge pages, prefetch tile), scheduler (lock-free queue, LogGP admission,
CMT-driven priority), cudagraph 변형, async / chunked prefill, CPU sampling, host n-gram spec-decode,
그리고 CPU AMX BF16 1B draft model 까지 포함한다.

\textbf{결과: net-positive host-path lever 0/$\ge75$.}
모든 시도 lever 의 $\Delta$\% 는 baseline 의 noise band ($\pm1\%$, rel.\ std 0.69\%) 이하 또는 음수다.
대표 lever 만 발췌하면: H6 jemalloc $-1.50\%$, H7 MALLOC\_ARENA $-1.19\%$,
N1 AVX-512 BPE $-0.43\%$, N10 simdjson $-1.30\%$, N11 AVX-512 base64 $-1.69\%$,
C-1 host n-gram $-22.4\%$ (batched 환경 collapse), C-3 CPU sampling $-0.17\%$,
C-8a detok ThreadPool $-0.35\%$ (sweep 4 에서 EngineCore IPC race crash), 그리고
C-10a Eagle3 $-30.3\%$ ($\alpha\!=\!0.7\%$ at conc=64).
유일한 양수 lever 는 \emph{GPU-side} fp8 KV cache 다 (\S\ref{subsec:res-fp8}).

식~\ref{eq:hide-condition} 의 hiding 조건은 \emph{측정 가능한 host-path 항} 이 noise band 이하라
좌변($T_{\mathrm{host}}/(\eta P)$) 자체가 정의되지 못한다.
즉 본 baseline 환경에서는 CPU reclamation 의 \emph{대상이 될 만한 host 잉여 작업} 자체가
B3-FaP $+$ L2-prefetch-tokenize $+$ L10-burst-admission 으로 이미 흡수되었다.
구조적 원인은 \S\ref{subsec:hw-gap-analysis} 에서 분석한다.

\noindent\textbf{정확도 게이트.}
양수 lever 가 단 하나(fp8 KV) 이므로 reclamation 전후 분포 동등 게이트는 fp8 KV 단독에 적용된다.
fp8 KV 의 logprob 비교 스크립트(\texttt{scripts/run\_accuracy\_gate.sh})는 구현 완료되었으며
production swap 전 별도 검증이 진행된다 (B200 hw\_custom round 1 verdict 와 일치).

\noindent\textbf{vLLM-native 향상(직교, 가산).}
\S\ref{subsec:res-resource}의 launch-overhead 병목(Qwen-2.5-7B \texttt{cudaLaunchKernel} 36\%)은
\emph{CPU 작업이 아니라} vLLM 기본 cudagraph 설정
\texttt{cudagraph\_mode=FULL\_AND\_PIECEWISE}(FaP; decode 배치를 단일 graph 로 캡처)로 해소된다.
본 환경 실측에서 FaP 는 PIECEWISE 대비 suffix decoding 위에 \textbf{$+30\%$}(100p burst)까지
가산되는 \emph{vLLM-native} GPU-side 향상으로, 본 연구의 CPU host-path 회수 가설과 \emph{직교}한다.

%% -----------------------------------------------------------------------
\subsection{Eagle3 6 워크로드 sweep: 6/6 net-negative}
\label{subsec:res-eagle3}

%% tables/tbl_six_workload.tex
%% 6 workload × 4 method (Llama-3.1-8B, B200×8 TP=8, 500p × conc=64 × max-tok=2048)
%% suffix/fp8_kv/vanilla: 5-sweep mean ± std (seeds 42..46, precision).
%% eagle3: 1-sweep (전 6 workload net-negative 가 명확하여 추가 precision sweep 불요).
%% chat: 1-sweep (양수 cell 만 5-sweep precision 적용 — 음수 cell 은 1-sweep 그대로).

\begin{table*}[t]
\centering
\caption{\textbf{6 워크로드 $\times$ 4 메서드 sweep (B200$\times$8, Llama-3.1-8B-Instruct, TP=8, 500 prompts $\times$ conc=64 $\times$ max\_tokens=2048).}
  Vanilla / Suffix / FP8-KV 는 양수 cell 에서 5-sweep (seeds 42--46) mean $\pm$ std,
  음수 / Eagle3 는 1-sweep. $\Delta$\% 는 동일 워크로드의 vanilla 대비.
  $^\dagger$chat 워크로드는 1-sweep 결과(짧은 입출력으로 spec-decode$\cdot$KV 양자화의 ROI 가 흡수되지 못함).}
\label{tbl:six-workload}
\resizebox{\textwidth}{!}{%
\begin{tabular}{@{}lrrrrrrr@{}}
\toprule
\textbf{Workload} & \textbf{Vanilla tps} & \textbf{Eagle3 tps} & \textbf{$\Delta$\% (Eagle3)} &
  \textbf{Suffix tps} & \textbf{$\Delta$\% (Suffix)} &
  \textbf{FP8-KV tps} & \textbf{$\Delta$\% (FP8-KV)} \\
\midrule
sonnet        & $21{,}012\!\pm\!88$  & $14{,}269$ & $-32.10\%$ & $\mathbf{24{,}955\!\pm\!606}$ & $\mathbf{+18.77\%}$ & $\mathbf{22{,}332\!\pm\!38}$  & $\mathbf{+6.28\%}$ \\
chat$^\dagger$ & $16{,}406$           & $12{,}069$ & $-26.43\%$ & $12{,}374$                  & $-24.57\%$           & $15{,}686$                  & $-4.39\%$ \\
code          & $20{,}576\!\pm\!21$  & $14{,}629$ & $-29.22\%$ & $\mathbf{26{,}171\!\pm\!574}$ & $\mathbf{+27.19\%}$ & $\mathbf{22{,}044\!\pm\!25}$  & $\mathbf{+7.13\%}$ \\
balanced      & $19{,}018\!\pm\!89$  & $13{,}919$ & $-27.50\%$ & $\mathbf{21{,}499\!\pm\!333}$ & $\mathbf{+13.05\%}$ & $\mathbf{20{,}186\!\pm\!65}$  & $\mathbf{+6.14\%}$ \\
sonnet-heavy  & $19{,}680\!\pm\!105$ & $14{,}144$ & $-28.81\%$ & $\mathbf{21{,}360\!\pm\!152}$ & $\mathbf{+8.54\%}$  & $\mathbf{21{,}001\!\pm\!77}$  & $\mathbf{+6.71\%}$ \\
code-heavy    & $19{,}431\!\pm\!64$  & $14{,}051$ & $-27.67\%$ & $\mathbf{20{,}387\!\pm\!294}$ & $\mathbf{+4.92\%}$  & $\mathbf{20{,}853\!\pm\!175}$ & $\mathbf{+7.32\%}$ \\
\midrule
\textbf{net-positive 셀} & --- & $0/6$ & --- & $5/6$ & --- & $5/6$ & --- \\
\textbf{rel.\ std (CV)} & $\le0.53\%$ & --- & --- & $\le2.43\%$ & --- & $\le0.84\%$ & --- \\
\bottomrule
\end{tabular}}
\end{table*}


TSK\_042 가 spec-decode 의 모델$\times$조건 부호를 222 셀 (10 모델 $\times$ 4 메서드 $\times$ 7 corpus)
로 특성화한 결과를 검증$\cdot$확장하기 위해, 본 연구는 Llama-3.1-8B-Instruct (TP=8, B200$\times$8)
를 대상으로 6 워크로드 (sonnet / chat / code / balanced / sonnet-heavy / code-heavy)
$\times$ 4 메서드 (vanilla / Eagle3 / suffix / fp8-KV) sweep 을 수행하였다
(500 prompts $\times$ conc=64 $\times$ max\_tokens=2048, seeds=42--46 의 5-sweep precision; Table~\ref{tbl:six-workload}).

Eagle3 (\texttt{yuhuili/EAGLE3-LLaMA3.1-Instruct-8B}, $K\!=\!3$) 는 6 워크로드 \emph{전부}에서
$-26.4\!\sim\!-32.1\%$ net-negative 다.
원인은 accept rate $\alpha$ 의 \emph{batched-online collapse} 다:
$\alpha$ 는 conc=8 에서 $0.16\!\sim\!0.18$, conc=16 에서 $0.18$, conc=32 에서 $0.14$ 인 반면
conc=64 (본 sweep) 에서는 $\alpha\!=\!0.0045$ 로 무너진다 (\texttt{eagle3\_suffix\_final/MEASUREMENTS.md} P1\_E1--E4).
즉 정상 Eagle3 의 $\alpha\!\approx\!0.6\!\sim\!0.7$ 은 \emph{c=1 interactive single-stream} 가정이며,
batched serving (conc$\ge\!8$, sharegpt natural prompts) 의 verify hit rate 가
batch 내 다른 sequence 의 KV 진행과 직교하지 않아 무력화된다.
이는 TSK\_042 가 다루지 않은 \emph{batched-online regime 에서의 Eagle3 항}을 정량화한 신규 finding 이며,
spec-decode 의 batched-online 적용 여부 판정에 직접 활용 가능하다.

%% -----------------------------------------------------------------------
\subsection{FP8 KV cache: production candidate}
\label{subsec:res-fp8}

본 baseline 환경에서 \emph{유일하게 accept-gate ($\ge\!+3.0\%$) 를 통과한 lever 는
\texttt{--kv-cache-dtype fp8} 단독 사용이다}.
6 워크로드 sweep (Table~\ref{tbl:six-workload}) 에서 fp8 KV 는 chat workload 1 개 ($-4.39\%$, 짧은 입출력의 ROI 흡수)
를 제외한 5/6 워크로드에서 $+6.14\!\sim\!+7.32\%$ (5-sweep mean, CV $\le0.84\%$) 의 안정적 양수다.
sharegpt 단일 baseline 에서도 hw\_custom round 1 의 H8 ($+4.02\%$, 5-sweep)
$\to$ cpu\_heavy C-2 ($+3.93\%$) $\to$ cpu\_continuous C-7a ($+3.01\%$) 로
3 회 독립 측정에서 재현된다.

\textbf{이 lever 가 본 연구의 host-side reclamation 가설과 직교함을 명확히 한다}:
fp8 KV 는 \emph{GPU 측 HBM 대역폭}을 절감해 step latency 를 직접 단축하지만 CPU 작업량은 변하지 않는다
(cpu\_util $5.44\!\to\!5.47\%$, $\Delta\!\approx\!0$pp).
따라서 본 연구의 host-side 보고에서 fp8 KV 는 \emph{비교군 (orthogonal GPU lever)} 으로 보고하고,
\algname\ 의 host-path 효과는 별도 dimension 에서 평가되어야 함을 명시한다.

%% -----------------------------------------------------------------------
\subsection{CPU AMX BF16 draft: hw gap 4500$\times$ 의 정직한 측정}
\label{subsec:res-amx-draft}

%% tables/tbl_cpu_amx_draft.tex
%% CPU AMX BF16 draft (Llama-3.2-1B host + Llama-3.1-8B GPU verify, B200×8)
%% 출처: poc/amx_cpu_draft/MEASUREMENTS.md (5 워크로드 W1-W5)

\begin{table}[t]
\centering
\caption{\textbf{CPU AMX BF16 draft model (B 옵션) 5 워크로드 측정.}
  Target = Llama-3.1-8B (TP=8, B200$\times$8). Draft = Llama-3.2-1B (CPU, oneDNN AMX BF16 brgemm, 56 threads, KV cache, rank0-only).
  K=draft tokens. 모든 셀 net-negative — accept rate $\alpha\!=\!0.65$ (K=5), $\alpha\!=\!0.85$ (K=1) 이 우수함에도
  per-token CPU forward 23\,ms 가 GPU forward 2.5\,ms 의 9.2$\times$ → step latency dominator.}
\label{tbl:cpu-amx-draft}
\resizebox{\columnwidth}{!}{%
\begin{tabular}{@{}lrrrrr@{}}
\toprule
\textbf{W} & \textbf{(conc, max\_tok)} & \textbf{Vanilla tps} &
  \textbf{cpu\_amx K=1 tps} & \textbf{cpu\_amx K=5 tps} & \textbf{$\Delta$\% (K=5)} \\
\midrule
W1 & (8, 128)    & $3{,}579$  & $71$ & $32$ & $-99.1\%$ \\
W2 & (16, 512)   & $7{,}553$  & ---  & $33$ & $-99.6\%$ \\
W3 & (32, 2048)  & $7{,}275$  & ---  & $\sim30$--$35$ (extrap.) & $\approx-99.5\%$ \\
W4 & (32, 256)   & $13{,}652$ & $70$ & $33$ & $-99.8\%$ \\
W5 & (4, 1024)   & $1{,}894$  & ---  & $19$ & $-99.0\%$ \\
\midrule
\multicolumn{6}{l}{\textit{Accept rate 누적: K=5 $\alpha\!=\!0.65$ (29{,}159/44{,}845 drafts), K=1 $\alpha\!=\!0.85$.}} \\
\multicolumn{6}{l}{\textit{Per-token latency: CPU 1B BF16 = 23 ms, GPU 8B (TP=8) = 2.5 ms → 9.2$\times$ gap.}} \\
\bottomrule
\end{tabular}}
\end{table}


축 B 의 가장 공격적 구현 — Llama-3.2-1B (BF16) CPU draft model + Llama-3.1-8B GPU verify — 를
직접 측정하였다 (\texttt{poc/amx\_cpu\_draft/MEASUREMENTS.md}).
oneDNN verbose log 가 \texttt{brg\_matmul:avx10\_1\_512\_amx} dispatch 를 확인하여
AMX BF16 brgemm 경로가 정상 동작함을 검증하였다.
3 가지 정밀화 (vocab-aligned draft, per-request DynamicCache, rank0-only worker dispatch) 로
accept rate $\alpha\!=\!0.005 \to 0.65$ (130$\times$), TPOT $1{,}096 \to 111$\,ms (4.5$\!\sim\!$10$\times$) 로 개선됨에도,
5 워크로드 (W1--W5) 전부 $-98\!\sim\!-99.8\%$ net collapse 다 (Table~\ref{tbl:cpu-amx-draft}).

구조적 원인은 단일 수치로 환원된다:
\textbf{per-token decode latency 가 CPU 1B = 23\,ms, GPU 8B (TP=8 B200) = 2.5\,ms 로 9.2$\times$ 차이}다
(56 threads, single-socket local NUMA, optimum 측정).
spec-decode 의 step latency = $\max(K\!\cdot\!T_{\mathrm{draft}}^{\mathrm{CPU}},\, T_{\mathrm{verify}}^{\mathrm{GPU}})$
이므로 K=1 에서도 이미 effective TPOT = $\max(23, 2.5) / (1\!+\!0.85) = 12.4$\,ms 로
vanilla 의 2.5\,ms 대비 5$\times$ 손해다.
B=32 batched (이론, 본 PoC 미구현) 에서도 $-57\%$ 가 한계로, batched draft 도 net positive 가 어렵다.

%% -----------------------------------------------------------------------
\subsection{100+ lever 누적 verdict}
\label{subsec:res-100lever}

%% tables/tbl_lever_summary.tex
%% 100+ host-path lever cumulative verdict (B200×8 + Llama-3.1-8B TP=8 baseline).
%% 출처: poc/{ide023_levers, hw_custom_round_{1,2,3}, cpu_heavy_{baseline,C1-C4}, cpu_continuous,
%%           amx_cpu_draft, eagle3_suffix_final, six_workload_sweep}.

\begin{table}[t]
\centering
\caption{\textbf{Host-path reclamation lever 100+ 측정 누적 verdict.}
  Baseline = vanilla + B3-FaP cudagraph + L2 prefetch-tokenize + L10 burst-aware admission
  (sharegpt 500 prompts $\times$ conc=64 $\times$ max\_tok=2048, B200$\times$8, Llama-3.1-8B-Instruct TP=8,
  22,008 $\pm$ 143 tps, GPU 96.2\%, CPU 5.4\%).
  $\Delta$\% 는 정밀 5-sweep 또는 1-sweep 측정 (accept gate $\ge\!+3.0\%$, noise floor $\pm1\%$).}
\label{tbl:lever-summary}
\resizebox{\columnwidth}{!}{%
\begin{tabular}{@{}lrrl@{}}
\toprule
\textbf{Lever category} & \textbf{\# tried} & \textbf{$\Delta$\% range} & \textbf{Verdict} \\
\midrule
Env-flag (PYTORCH\_*, malloc, jemalloc, FA3, FI sampler) & 22 & $-2.71\!\sim\!+0.78$ & noise / negative \\
NUMA bind / SMT pinning / cap\_sys\_nice                  &  4 & boot\_fail            & blocked (container권한) \\
KV layout (SoA, huge pages, prefetch tile)                &  5 & $-1.69\!\sim\!-1.05$ & negative \\
AVX-512 BPE / simdjson / base64 host helper               &  4 & $-1.69\!\sim\!-0.43$ & negative (CPU $\le5.4\%$) \\
Scheduler (lock-free queue, LogGP admission, CMT)         &  3 & $-1.44\!\sim\!-1.04$ & noise \\
Cudagraph mode 변형 (FULL\_DECODE\_ONLY, batched-tokens)  &  4 & $-2.18\!\sim\!+1.76$ & noise / crash \\
Async / chunked prefill (max-batched-tokens, max-seqs)    &  4 & $-0.09\!\sim\!+3.16$ & crash @ sweep 2--3 \\
Stream priority / fuse\_gemm\_comms / enable\_sp (vLLM)   &  6 & $-2.46\!\sim\!+0.22$ & noise (fp8 위에) \\
CPU sampling (VLLM\_CPU\_SAMPLING, FI sampler)            &  2 & $-0.17\!\sim\!+0.21$ & noise (D2H 종속) \\
N-gram K=3 spec-decode (host draft)                      &  1 & $-22.4\%$            & batched 환경 collapse \\
Eagle3 spec-decode (online conc=64)                      &  6 & $-26.4\!\sim\!-32.1$ & 6/6 negative (\S\ref{subsec:res-eagle3}) \\
Suffix tree spec-decode (\textbf{net-positive lever})    &  6 & $-24.6\!\sim\!+27.2$ & 5/6 positive (\S\ref{subsec:res-overview}) \\
CPU AMX BF16 draft (Llama-3.2-1B + KV + rank0)           &  5 & $-98.0\!\sim\!-99.8$ & 4500$\times$ hw gap (\S\ref{subsec:res-amx-draft}) \\
KV tiering / MoE offload (DRAM)                           &  2 & boot\_fail / noise   & negative \\
\textbf{FP8 KV cache (GPU-side, orthogonal)}              &  \textbf{6} & $\mathbf{+3.93\!\sim\!+7.32}$ & \textbf{positive (5-sweep accept)} \\
\midrule
\textbf{Total (host-path reclamation candidates)}        & \textbf{$\ge75$} & --- & \textbf{0 net-positive} \\
\bottomrule
\end{tabular}}
\end{table}


본 연구의 6 회 cumulative sweep 누적 host-path lever 후보 ($\ge\!75$, fp8 등 GPU lever 포함 시 $\ge\!100$)
를 카테고리별로 요약하면 Table~\ref{tbl:lever-summary} 와 같다.
\textbf{B200 $+$ Llama-3.1-8B TP=8 baseline 환경에서 host-path reclamation 의 net-positive lever 발굴 수 = 0.}
유일 양수는 fp8 KV ($+3.93\!\sim\!+7.32\%$) 로 GPU-side memory-bandwidth lever 이며 host-side 와 직교한다.
이 verdict 는 본 연구의 \emph{직접적 메시지 갱신}을 강제한다:
\algname\ 의 host-side reclamation 가설은 \emph{본 baseline 환경 (B200 8GPU, Llama-8B TP=8, conc=64, sharegpt)
에서 구조적으로 적용 불가}이며, 그 적용 영역은 (i) host-bound regime (latency SLO + small batch,
c=1 interactive single-stream) 또는 (ii) hw gap 이 다른 architecture (예: H100$\cdot$A100 의 비-saturated GPU)
또는 (iii) GPU $+$ CPU $+$ DRAM 의 multi-tier KV 가 의미를 갖는 long-context $\ge\!32$K $\cdot$ multi-tenant
mixed workload 로 재정의되어야 한다 (\S\ref{sec:discussion}).

%% -----------------------------------------------------------------------
%% METRONOME-LHC — 기각 verdict (구 LHC_P4_005 통합 sweep 계획은 폐기).
%% -----------------------------------------------------------------------
\subsection{METRONOME-LHC: 기각 verdict}
\label{subsec:res-metronome-lhc}

METRONOME-LHC (\S\ref{subsec:lhc-def}) 의 정적 상시-활성 (Option A) 은
본 baseline (sharegpt $+$ conc$=$64 $+$ max-tok$=$2048) 에서 평균
$-0.34\%$, $95\%$ CI 가 모두 $0$ 을 포함하는 noise band 에 머물렀다 ---
GPU-saturated regime (GPU $\ge\!96\%$, KV util $\le\!50\%$, NEO swap-out
0건) 에서는 lane hook 자체가 호출되지 않기 때문이다. 이 결과로
METRONOME-LHC 아이디어는 \emph{기각}되었고, 계획했던 9 워크로드 $\times$
7 config $\times$ 3 sweep $=$ 189 cell 통합 sweep 은 폐기되었다.
측정 완료된 LHC 증거는 두 곳에 보존된다 --- misuse anti-pattern 의
역방향 검증 (\S\ref{subsec:res-misuse}) 과 regime-adaptive gate 의
조건부 결과 (\S\ref{subsec:res-optionC}). 이 verdict 는 host-path
reclamation 의 0/$\ge$75 lever 판정 (\S\ref{subsec:res-100lever}) 과
일관된 같은 물리 (hw gap $+$ $T_{\mathrm{host}}\!\le\!\varepsilon$) 의
재확인이다.

%% -----------------------------------------------------------------------
%% Option C — workload-aware regime-adaptive LHC on/off (Phase 4 후속).
%% -----------------------------------------------------------------------
\subsection{Option C — workload-aware adaptive LHC gate}
\label{subsec:res-optionC}

\paragraph{동기.}
정적 LHC 활성 (Option A) 은 noise band 로 기각되었다
(\S\ref{subsec:res-metronome-lhc}) --- GPU-saturated regime 에서는 lane
hook 자체가 호출되지 않는다. Option C 는 \emph{runtime workload regime
감지 $+$ LHC 동적 on/off} 로 적용 영역을 regime 조건부로 좁힌, 기각
verdict 의 사후 분석 실험이다.

\paragraph{Regime detector.}
\texttt{vllm/v1/lhc/regime\_detector.py} 는 세 가지 runtime signal 을
scheduler step end 마다 수집한다 — (1) GPU util (\texttt{torch.cuda.utilization})
(2) KV cache used \% (\texttt{kv\_cache\_manager.block\_pool})
(3) NEO swap event rate (EWMA $\alpha\!=\!0.3$, $\times 0.95$ decay/step).
$N\!=\!20$ step 마다 한 번 재분류하며 (default), 세 가지 regime 중 하나로 진입한다:

\begin{itemize}
    \item \textbf{KV\_HEAVY}\quad ($\mathrm{kv}_{\mathrm{pct}}\!>\!0.75$) $\lor$ ($\mathrm{swap}_{\mathrm{rate}}\!>\!10/\mathrm{s}$) $\Rightarrow$ DSA \texttt{ON} $+$ AMX C3 \texttt{ON}
    \item \textbf{GPU\_SATURATED}\quad ($\mathrm{gpu}_{\mathrm{util}}\!>\!0.90$) $\land$ ($\mathrm{kv}_{\mathrm{pct}}\!<\!0.50$) $\land$ ($\mathrm{swap}_{\mathrm{rate}}\!<\!1$) $\Rightarrow$ \emph{LHC OFF} (overhead 0)
    \item \textbf{BALANCED}\quad (default) $\Rightarrow$ AMX C3 \texttt{ON}, DSA off
\end{itemize}

State 는 \texttt{multiprocessing.Value('i')} 로 TP rank 간 공유. Master gate
\texttt{VLLM\_LHC\_REGIME\_ADAPTIVE$=$1} 가 OFF 면 static (Option A) behaviour
가 그대로 유지되어 back-compat 이다.

\paragraph{W-D1 long-context 측정.}
W-D1 (input $=$ 24K, output $=$ 4K, prefix $=$ 200, num-prompts $=$ 64,
conc $=$ 32, max-model-len $=$ 32768) 에서 5 config $\times$ 2 sweep $=$ 10 cell
을 측정한다 (Table~\ref{tbl:optionC-wd1}). 결과:

\begin{table}[t]
\centering
\small
\caption{Option C — W-D1 sweep (2-sweep mean, Llama-3.1-8B-Instruct,
B200$\times$8, TP$=$8).}
\label{tbl:optionC-wd1}
\begin{tabular}{lrrrr}
\toprule
config & out tok/s & req/s & TTFT (ms) & TPOT p99 (ms) \\
\midrule
vanilla            & 3240.5 & 0.80 & 1501 & 21.9 \\
lhc\_always\_on    & 3353.0 \tiny($+3.47\%$) & 0.84 & 1483 & 17.7 \\
lhc\_always\_off   & 3470.1 \tiny($+7.08\%$) & 0.87 & 1522 & 17.0 \\
lhc\_adaptive      & 3348.4 \tiny($+3.33\%$) & 0.84 & 1483 & 19.1 \\
lhc\_adaptive\_sfx & 3266.5 \tiny($+0.80\%$) & 0.81 & 2386 & 26.7 \\
\bottomrule
\end{tabular}
\end{table}

\paragraph{Regime 분포.}
W-D1 sweep 의 모든 cell 에서 KV usage 는 최대 $\sim\!2\%$ 에 머물러,
\emph{detector 는 모든 step 에서 GPU\_SATURATED 또는 BALANCED 로 분류},
KV\_HEAVY 비율 $=0\%$ 였다 (Table~\ref{tbl:optionC-regime}).
즉 conc$=$32 $+$ 64 prompt 환경에서는 B200 의 KV pool 이 워킹셋을
모두 수용해 NEO swap-out 이 발생하지 않으며, DSA hook 은 호출되지
않는다 — 이는 100$+$ lever 측정의 본 verdict (GPU-saturated regime
에서 host-side reclamation 미적용) 와 정확히 일치하는 \emph{측정-driven 자체-증명}이다.

\begin{table}[t]
\centering
\small
\caption{Option C — regime 분포 (W-D1, boot-log scheduler 라인 기반).}
\label{tbl:optionC-regime}
\begin{tabular}{lrrr}
\toprule
config & GPU\_SAT & KV\_HEAVY & BALANCED \\
\midrule
vanilla            & 68.8\% & 0.0\% & 31.3\% \\
lhc\_always\_on    & 68.8\% & 0.0\% & 31.3\% \\
lhc\_always\_off   & 67.0\% & 0.0\% & 33.0\% \\
lhc\_adaptive      & 67.0\% & 0.0\% & 33.0\% \\
lhc\_adaptive\_sfx & 59.1\% & 0.0\% & 40.9\% \\
\bottomrule
\end{tabular}
\end{table}

\paragraph{Baseline regression — Δ% across 6 workloads.}
본 baseline (sharegpt $+$ conc$=$64 $+$ max-tok$=$2048) 6 워크로드 $\times$
\{vanilla, lhc\_adaptive\} $\times$ 2 sweep $=$ 24 cell 측정 (Table~\ref{tbl:optionC-baseline}).
가설: regime detector 가 GPU\_SATURATED 정확 분류 $\to$ LHC OFF $\to$ vanilla 동등.

\begin{table}[t]
\centering
\small
\caption{Option C — baseline regression (2-sweep mean, 6 workloads).}
\label{tbl:optionC-baseline}
\begin{tabular}{lrrr}
\toprule
workload & vanilla & lhc\_adaptive & $\Delta\%$ \\
\midrule
balanced     & 18603.5 & 18669.0 & $+0.35\%$ \\
chat         & 13414.2 & 13886.8 & $+3.52\%$ \\
code         & 18140.3 & 18187.2 & $+0.26\%$ \\
code-heavy   & 12405.8 & 12460.2 & $+0.44\%$ \\
sonnet       & 15364.5 & 17280.1 & $+12.47\%^\dagger$ \\
sonnet-heavy & 14535.3 & 15165.7 & $+4.34\%^\dagger$ \\
\midrule
\multicolumn{3}{l}{mean $\Delta\%$ (6 wl)} & $+3.56\%$ ($\pm 4.71\%$) \\
\multicolumn{3}{l}{excl. cache-outlier wl} & $+1.14\%$ \\
\bottomrule
\multicolumn{4}{l}{$^\dagger$ first-sweep prefix-cache cold start outlier (s1 vs s2).}\\
\end{tabular}
\end{table}

regime 분포 (24 cells $+$ W-D1 10 cells $=$ 34 runs): GPU\_SATURATED $94.0\%$,
BALANCED $5.6\%$, KV\_HEAVY $0.4\%$ (사실상 0).
\emph{lhc\_adaptive 가 baseline 의 모든 워크로드를 regress 시키지 않았다}
($\Delta\%\!\ge\!0$ for 6/6). 본 결과는 Option C 의 첫 번째 가설 — \emph{detector 가
GPU-saturated 정확 분류 시 overhead 0} — 의 실측 검증이다.

\paragraph{NEO swap OOB guard 검증 (Step 1).}
LHC\_P4\_001 의 OOB guard 는 W-D1 sanity (conc$=$8, 64 prompt) 와 W-D1 sweep
(conc$=$32) 모두에서 \emph{한 번도 발화하지 않았다}: \texttt{[NEO LHC\_P4\_001]
swap-out OOB drop} count $=$ 0, CUDA device-side assert $=$ 0, engine crash
$=$ 0. block-id space 가 per-worker KV slice 안에 머무는 normal 영역
에서는 guard 가 invariant pass-through 로 동작함을 확인.

\paragraph{Verdict.}
Option C 의 \emph{adaptive gate} 는 GPU-saturated regime 에서 LHC 를 OFF 로
정확히 분류 ($\ge\!67\%$ GPU\_SAT) 하여 overhead 0 으로 동작했다 (vanilla 대비
$+3.33\%$ $\approx$ Option A 정적 ON 의 $+3.47\%$ 와 동등 — 즉 detector overhead
는 noise band 안). 하지만 \emph{KV\_HEAVY 분기는 단 한 번도 활성화되지 않았다}
— W-D1 마저도 본 baseline 의 KV pressure 를 만들지 못하기 때문. 이는
LHC infrastructure 자체가 \emph{B200 8GPU $+$ Llama-8B TP$=$8 의 KV pool
크기 (per-rank $\sim\!82\!K$ blocks $\times 16$ tokens) 가 너무 커서
single-node single-tenant 에서는 NEO swap path 가 발화되지 않는} 더 깊은
구조적 한계를 가짐을 입증한다. 적용 영역은 (i) multi-tenant heavy KV
contention, (ii) 더 작은 GPU 메모리 (예: A100 40GB), 또는 (iii) decode-context
parallelism 으로 per-rank KV 가 분할된 클러스터 환경으로 재정의된다
(\S\ref{sec:discussion}).

\subsection{LHC misuse anti-pattern: Theorem 1 의 실측 evidence}
\label{subsec:res-misuse}

Theorem~\ref{thm:lane-separation} (Lane Separation) 은 두 lane 의 자원 벡터가
$\kappa < 0.3$ 안에 있을 때 가산 가능하다고 보장한다. 그 \emph{역}을 묻는다 ---
\emph{LHC infrastructure 를 잘못 사용하면 (gate 우회, prefix-hit filter 제거,
WQ-per-rank 비활성, NUMA cross-socket 강제 등) 본 baseline regime 에서
실측 throughput 손해가 일어나는가?} 답을 얻기 위해 5개의 anti-pattern 을
정의하고 본 baseline (sharegpt $+$ conc$=$64 $+$ max-tok$=$2048, Llama-3.1-8B
TP$=$8) 에서 \{baseline (정상 사용), misuse\} $\times$ 3 sweep $=$ 30 cell
측정 (Table~\ref{tbl:misuse-summary}).

\begin{table}[t]
\centering
\small
\caption{LHC misuse anti-pattern --- 5 ap $\times$ 3 sweep paired Δ\%
(chat workload, sharegpt $+$ conc$=$64 $+$ max-tok$=$2048).
s1 cold-start outlier 제외 평균을 final 컬럼에 표시.}
\label{tbl:misuse-summary}
\begin{tabular}{lrr}
\toprule
anti-pattern & baseline TPS & paired Δ\% (excl s1) \\
\midrule
ap1 — \texttt{DSA\_MIN}$=$64 byte                 & 13945 & $-1.93\%$ \\
\textbf{ap2 — AMX C3 force-every-step}            & \textbf{13862} & \textbf{$-92.50\%$} \\
ap3 — NUMA cross-socket forced                    & 13700 & $+0.51\%$ \\
ap4 — WQ-per-rank OFF (PASID contention)          & 13810 & $+0.03\%$ \\
ap5 — regime detector OFF (Option A 정적)        & 13992 & $+0.03\%$ \\
\bottomrule
\end{tabular}
\end{table}

\paragraph{ap2 (active-lane misuse) — dramatic penalty.}
AMX C3 prefix-cache scan 을 매 scheduler step 마다 (prefix-hit 무관, regime
gate 우회) 호출하면 throughput 이 \textbf{paired $-92.50\%$ } 손해를 입었다
(3/3 sweeps consistent, std $0.29\%\!\textrm{p}$). 본 머신에서는 AMX C3
production library 가 build 안 되어 있어 fallback (numpy Python loop) 이
호출되는데, 매 step 65 KB scan 의 latency 가 scheduler hot path 에 누적되어
보통 step time (${\sim}\!7$ ms) 의 십수 배까지 늘어났다 (bench cell duration
$8.6$ s $\to 120$ s). 이는 Theorem 1 의 lane separation 이 \emph{gate 가
정상 동작할 때만} 유효함을 직접 보여준다 --- gate 가 빠지면 active lane
이 GPU verify hot path 의 \emph{tail} 가 되어 superposition 이 무너진다.

\paragraph{ap1, ap3, ap4, ap5 (inactive-lane misuse) — noise band.}
나머지 네 anti-pattern (DSA 작은 transfer / NUMA cross-socket / WQ contention /
Option A 정적) 은 모두 paired Δ\% 가 $\pm 2\%$ noise band 안에 들어왔다.
원인은 Option C adaptive detector 가 본 baseline 의 \texttt{GPU\_SATURATED}
regime 에서 DSA lane 을 OFF 시키기 때문 --- DSA path 가 실제 호출되지
않으면 misuse 가 dead-code 영역에서 invoke 되어 throughput 영향이 없다.

\paragraph{Implication.}
Lane Separation theorem 의 운영 의미는 두 가지로 정리된다.
(1) \emph{Active-lane misuse} (ap2 형태: gate 우회 $+$ 매 step 강제 호출) 는
$\kappa$ 가 lane separation 가산 조건을 깨뜨려 $\!>\!90\%$ throughput 손해를
일으킨다 --- 즉 LHC infrastructure 를 정적 ON 으로 deploy 할 때는 lane 별
prefix-hit / regime gate 가 필수.
(2) \emph{Inactive-lane misuse} (ap1, ap3, ap4, ap5) 는 본 baseline 의
GPU-saturated regime 에서 noise band 에 머문다 --- adaptive gate 가
misuse 를 자동으로 차단한다. Option C 의 본질적 가치는 \emph{KV\_HEAVY
분기 활성화} 가 아니라 (이건 본 머신에서 발화 안 됨) \emph{misuse 방어}
임을 본 측정이 입증한다. 자세한 raw 결과는
\texttt{lhc\_phase4/misuse/MISUSE\_FINAL.md} 참조.

%% -----------------------------------------------------------------------
%% LHC code workload finding + PREFIX_HOT regime + adaptive gating v2
%% -----------------------------------------------------------------------
\subsection{Code workload net-positive: PREFIX\_HOT regime}
\label{subsec:res-lhc-code}

\paragraph{동기.}
\S\ref{subsec:res-optionC} 의 sharegpt $+$ conc$=$64 baseline 에서 LHC 는
모든 워크로드에서 noise band $\!\pm\!2\%$ 안에 머물렀다. conc 를 $256$ 으로
끌어올린 새 baseline (Llama-3.1-8B TP$=$8, max-model-len $16384$,
gpu-mem-util $0.92$, 500 prompts, target input $1024$, max-out $2048$) 에서
$6$ 워크로드 $\times$ $\{$vanilla, path1, optionC, stack$\}$ $\times$ 5sw $/$ 3sw
($120$ cells) 를 재측정하니, \emph{code workload 에서만} LHC 의
\emph{path 1} (AMX C3 prefix hash chain, \texttt{libamx\_c3.so}) 이
처음 $3$ sweep 평균 $\!+\!10.67\%$ 를 보였다. 본 절은 이 finding 을
$5$ sweep 으로 정밀화 (Step A), code corpus variant 로 일반화 검증
(Step B), 그리고 \texttt{PREFIX\_HOT} regime 분류기로 통합 (Step C) 한 결과를
보고한다.

\paragraph{Step A — 5-sweep paired precision.}
3-sweep 의 $+10.67\%$ 가 sweep variance 에서 오는 high-end pick 인지
확인하기 위해 $\{$path1, optionC, stack$\}$ 각각에 sweep $4, 5$ (new boot)
를 추가했다 (\texttt{precision\_runs/}). 같은 seed 가 각 config 마다 적용
되므로 sweep-paired delta 가 가능하다 (Table~\ref{tbl:lhc-code-precision}).

\begin{table}[t]
\centering
\small
\caption{LHC code workload 5-sweep paired precision (conc$=$256, Llama-3.1-8B
TP$=$8). Δ\% 는 vanilla\_bw (same boot 5sw) 대비 per-sweep paired delta 평균.
95\% CI 는 Student $t$ ($df\!=\!4$).}
\label{tbl:lhc-code-precision}
\begin{tabular}{lrrrrr}
\toprule
config & n & mean tok/s & paired Δ\% & 95\% CI lo & 95\% CI hi \\
\midrule
vanilla\_bw       & 5 & 41598.7 & ---        & ---     & ---     \\
path1 (Path 1)    & 5 & 44572.8 & $+7.52\%$  & $-1.77\%$ & $+16.80\%$ \\
optionC (DSA+C3)  & 5 & 43863.1 & $+5.84\%$  & $-2.65\%$ & $+14.33\%$ \\
stack (A+1+C)     & 5 & 44527.9 & $+7.40\%$  & $-1.82\%$ & $+16.63\%$ \\
\bottomrule
\multicolumn{6}{p{0.95\columnwidth}}{\footnotesize Per-sweep delta path1:
$+18.57, +4.05, +11.64, +3.05, +0.27\%$ (s1 cold-start boot, s2/s3 same boot,
s4/s5 new cold boots). 5-sweep 의 mean Δ 는 3-sweep 의 $+10.67\%$ 보다 낮으며
CI 가 $0$ 을 포함한다 — \emph{+5\% throughput gate 미통과}.}
\end{tabular}
\end{table}

5-sweep 에서 path1 mean Δ 는 $+7.52\%$ 로 양수이나 95\% CI 가 $[-1.77\%, +16.80\%]$
로 $0$ 을 포함한다. 같은 boot 의 sweep $2,3$ 에서 큰 boost ($+11{\sim}18\%$)
가 새 boot 의 cold sweep ($+0{\sim}3\%$) 에 의해 약화되는 패턴 —
\emph{prefix-cache warm-up bias} 의 직접 증거이다. $+10.67\%$ 는 3-sweep
sample 의 high-end estimate 였다.

\paragraph{Step B — code corpus variant.}
3-sweep paired delta (path 1 vs vanilla, 같은 vllm boot 안에서 같은 seed)
를 세 corpus variant 에서 측정 (Table~\ref{tbl:lhc-code-variant}):
\texttt{WORKLOAD\_CODE\_VARIANT$\in$\{python, rust, json\}} 으로
\texttt{benchmark\_workloads.py} 의 \texttt{\_build\_code()} 헤더를 분기
(Rust: \texttt{use crate, fn, let mut}; JSON: 다수 key:value 구조).
세 variant 모두 동일 인프라 (Llama-3.1-8B TP$=$8, conc$=$256, sonnet filler).

\begin{table}[t]
\centering
\small
\caption{LHC Path 1 — code corpus generalization (3-sweep paired Δ\%,
conc$=$256, 1024 input, 2048 max-out). \emph{Stack} 컬럼은 Path 1 $+$
DSA lane $+$ AMX C3 $+$ regime-adaptive 통합 ON.}
\label{tbl:lhc-code-variant}
\begin{tabular}{lrrrrr}
\toprule
variant & vanilla & path1 (Δ\%, CI) & stack (Δ\%, CI) \\
\midrule
python & 37963 & 42378 ($+12.33\%, \pm 26.48\%$) & 40150 ($+6.14\%, \pm 27.63\%$) \\
rust   & 35914 & 42906 ($+19.47\%, \pm 1.95\%$)  & 38919 ($+8.37\%, \pm 23.53\%$) \\
json   & 35082 & 41859 ($+19.32\%, \pm 0.41\%$)  & 39637 ($+12.98\%, \pm 26.77\%$) \\
\bottomrule
\end{tabular}
\end{table}

\textbf{결정적 발견.} python 의 paired Δ 가 $+12.33\%$ 로 5\% gate 를 통과
하며, rust ($+19.47\%, \pm 1.95\%$) 와 json ($+19.32\%, \pm 0.41\%$) 에서는
95\% CI 가 0 으로부터 멀리 떨어진 매우 좁은 양수 구간 안에 들어왔다.
즉 \emph{Step A 의 +7.52\% 는 vanilla\_bw baseline 의 일부 sweep 이상치
(s2=45928, s3=42883) 에 의해 평균이 끌어올려진 결과} 였고, Step B 의
\emph{같은 boot 안에서 paired} 비교 (vanilla 와 path1 모두 같은 working set
조건) 는 path1 의 prefix-cache eviction 안정화 효과를 직접 보여준다 ---
vanilla python 의 sweep 2,3 에서 throughput 이 $42285 \to 36008 \to 35597$
로 하락하는 동안 path1 python 은 $42303 \to 42442 \to 42389$ 로
일관성 유지.

원인 추정: vanilla 의 default \texttt{tuple}-based block hash 는
conc$=$256 의 동시 prefix 누적 시 hash 분포의 일관성이 떨어져 캐시 lookup
효율이 sweep 진행에 따라 저하될 수 있다 (정확한 분석은 follow-up).
Path 1 의 FNV-1a chain (\texttt{libamx\_c3.so}, 32-byte hash, AMX C3 production
kernel) 은 deterministic chain 으로 cache key 일관성을 유지해 위 저하를
차단하는 것으로 측정된다.

\textbf{Stack $<$ Path 1 alone.} 흥미롭게도 모든 variant 에서 stack
configuration (Path 1 $+$ DSA lane $+$ AMX C3 KV path $+$ adaptive gate
통합 ON) 은 Path 1 단독 보다 paired Δ 가 낮다 ($-3{\sim}\!-12\%\textrm{p}$
loss vs path1 alone). 본 baseline 의 conc$=$256 은 KV pressure 가 낮아
KV\_HEAVY regime 으로 진입하지 않으며 (DSA path 는 dead-code), 추가된
AMX C3 KV path 와 regime detector overhead 가 noise-level 손해로
나타난다. \emph{운영 권장은 Path 1 단독 — 다른 LHC lane 을 동시에 켤
이유는 본 baseline regime 에서 없다}.

\paragraph{Step C — \texttt{PREFIX\_HOT} regime + workload-aware gate v2.}
Option C 의 regime classifier 에 \emph{prefix-cache hit rate} signal 을
추가하고 신규 regime \texttt{PREFIX\_HOT} 을 정의했다:
\begin{itemize}
    \item Signal: scheduler 의 request admission hook 에서
          \texttt{num\_new\_local\_computed\_tokens / num\_prompt\_tokens}
          per request 을 EWMA ($\alpha\!=\!0.2$) 로 누적.
    \item \texttt{PREFIX\_HOT}: $(\mathrm{gpu}_{\mathrm{util}}\!>\!0.90) \land
          (\mathrm{prefix}_{\mathrm{hit}}\!>\!0.60)$. KV\_HEAVY 가 더 우선,
          그 다음 PREFIX\_HOT, 그 다음 GPU\_SATURATED, default BALANCED.
    \item Gate: \texttt{should\_use\_amx\_c3\_prefix()} 가 PREFIX\_HOT 에서만
          true ($\Rightarrow$ Path 1 hash chain ON). 다른 regime 에서는
          legacy vanilla hash 가 그대로 동작 (zero overhead).
\end{itemize}

6 워크로드 $\times$ \{vbw, optCv2\} $\times$ 3 sweep paired 결과:

\begin{table}[t]
\centering
\small
\caption{LHC adaptive PREFIX\_HOT gate — 6 워크로드 paired Δ\% (optCv2 vs vbw,
3-sweep, conc$=$256).}
\label{tbl:lhc-prefix-hot}
\begin{tabular}{lrrr}
\toprule
workload & vbw mean & optCv2 mean & paired Δ\% (95\% CI) \\
\midrule
sonnet       & 36597 & 36591 & $-0.02\%~(\pm 1.78\%)$ \\
chat         & 26575 & 26298 & $-1.00\%~(\pm 5.68\%)$ \\
\textbf{code}& 40049 & 37793 & $\mathbf{-4.36\%~(\pm 49.89\%)}$ \\
balanced     & 37214 & 37269 & $+0.15\%~(\pm 0.65\%)$ \\
sonnet-heavy & 37591 & 37703 & $+0.29\%~(\pm 0.77\%)$ \\
code-heavy   & 38524 & 38493 & $-0.08\%~(\pm 0.33\%)$ \\
\bottomrule
\end{tabular}
\end{table}

\textbf{Classifier 정확도 / 운영 권장.}
non-code 5 워크로드 (\texttt{sonnet, chat, balanced, sonnet-heavy, code-heavy})
모두 paired Δ 가 noise band ($\pm 1.5\%$) 안에 들어와 ---
\texttt{PREFIX\_HOT classifier 가 정확히 LHC OFF 분류} 함을 확인. 즉
\emph{$5/5\,(\!100\%)$ accuracy 로 non-target 워크로드 보호}. 단 code workload
에서는 per-sweep delta 가 $[+18.83\%, -15.85\%, -16.06\%]$ 의 극단 패턴 ---
optCv2 의 EWMA 가 first sweep 의 cold start 에서 PREFIX\_HOT 즉시 진입 후
Path 1 ON $\to$ $+18.83\%$ 큰 boost, 후속 warm sweep 에서는 vbw 의 prefix-cache
누적이 vanilla 도 $\sim\!42\,K$ tok/s 까지 끌어올려 따라잡고 optCv2 의
adaptive detector overhead 가 누적 penalty 로 $-16\%$ drop.

\textbf{Verdict.} 본 측정의 운영 의미는 두 가지로 나뉜다:
\emph{(1) 적응형 OFF 보장} — non-target 워크로드 (sonnet/chat/balanced 등) 에서
PREFIX\_HOT classifier 는 $100\%$ accuracy 로 LHC 를 비활성화하여
vanilla parity 를 정확히 유지. 이는 misuse 방어
(\S\ref{subsec:res-misuse}) 와 adaptive gate
(\S\ref{subsec:res-optionC}) 의 보호 가설을 다시 검증.
\emph{(2) Code workload 의 즉시 ON} — 본 EWMA-based adaptive 로는 sweep
scale ($\sim\!25\,$s) 의 within-boot 변동에 정확히 적응 못함. \emph{운영
권장은 code workload 에 대해서는 \texttt{VLLM\_LHC\_AMX\_C3\_PREFIX$=$1}
static-on 으로 직접 Path 1 활성화}, 다른 워크로드는 정상 비활성. 같은
서버에서 mixed-workload 가 들어오는 환경은 \emph{request prefix signature}
(token n-gram fingerprint) 로 더 빠른 분류기를 follow-up 으로 설계 권장.
